% !TEX program = xelatex
\documentclass[11pt, a4paper]{article}

\usepackage{fontspec}
\setmainfont[
  Path=/usr/share/texmf/fonts/opentype/public/tex-gyre/,
  UprightFont=texgyrepagella-regular.otf,
  BoldFont=texgyrepagella-bold.otf,
  ItalicFont=texgyrepagella-italic.otf,
  BoldItalicFont=texgyrepagella-bolditalic.otf
]{TeX Gyre Pagella}

\usepackage{amsmath, amssymb, amsthm}
\usepackage[margin=1.25in]{geometry}
\usepackage{setspace}
\onehalfspacing
\usepackage{parskip}
\setlength{\parindent}{0pt}
\setlength{\parskip}{0.85em}
\usepackage[hidelinks]{hyperref}
\usepackage{titlesec}
\titleformat{\section}{\large\bfseries}{}{0em}{}[\vspace{0.2em}\hrule\vspace{0.4em}]
\titleformat{\subsection}{\normalsize\bfseries\itshape}{}{0em}{}
\usepackage{epigraph}
\setlength{\epigraphwidth}{0.62\textwidth}
\renewcommand{\epigraphflush}{center}

\theoremstyle{definition}
\newtheorem{definition}{Definition}
\newtheorem{theorem}{Theorem}
\newtheorem{proposition}{Proposition}
\newtheorem{corollary}{Corollary}
\theoremstyle{remark}
\newtheorem{remark}{Remark}

\usepackage{titling}
\setlength{\droptitle}{-2em}

\title{\textbf{Semantic Stability as a Repair Equilibrium}\\
{\large Threshold, Attenuation, and Irreversible Collapse\\
in Constraint-Bounded Semantic Trees}}
\author{Flyxion \\ \small Independent Researcher}
\date{2026}

\begin{document}

\maketitle

\begin{abstract}
A constraint-bounded semantic tree is a structure in which meaning propagates
from root to leaf subject to an attenuation process that weakens semantic
coherence with distance. We formalize such trees as repair equilibria in the
sense of the accompanying repair ontology: each node remains semantically
alive so long as the repair pressure reaching it --- the accumulated constraint
propagation from its ancestors --- meets or exceeds a local admissibility
threshold. The central result is the Semantic Stability Theorem, which gives
a branch coherence condition $C(\mathcal{B}) \geq 1$ for the persistence of
a connected semantic subtree. When this condition fails, the branch crosses
the boundary of its admissibility manifold and enters irreversible attenuation.
We prove that boundary nodes are the critical sites and that local boundary
failure propagates inward through the tree in a manner analogous to crack
propagation in materials under load. The RSVP field triple $(\Phi, v, S)$
provides a natural field-theoretic interpretation: semantic density, constraint
propagation velocity, and accessible future volume map onto the three field
variables, and irreversible attenuation corresponds to local entropy expansion
outrunning the repair flow of the vector field. The threshold is therefore not
an information threshold but a repair threshold. Semantic death is unrepaired
divergence.
\end{abstract}

\newpage

\epigraph{\textit{A word that is no longer repaired by use falls out of the
language. Not because it has been forgotten, but because the repair
infrastructure for it has collapsed.}}{}

\section{Introduction}

The companion essay establishes that objects and processes are emergent
descriptions of repair equilibria: configurations in which the rate of repair
meets or exceeds the rate of degradation for a sustained time horizon. That
framework applies at every scale from molecular biology to cosmology. This
paper applies it to a specific and formally tractable domain: the propagation
of semantic constraint through a hierarchical tree structure.

The question motivating the paper is: when does a semantic structure remain
coherent, and when does it collapse irreversibly? The question is not about
information in the Shannon sense --- bits stored, transmitted, received. It is
about constraint: does the semantic content at the root of a tree still exert
meaningful influence on interpretation at the leaves? And if that influence
is weakening --- through distance, through conceptual branching, through
competing interpretive pressures --- at what point does the weakening become
irreversible?

The repair framework gives a precise answer. A semantic tree persists as a
coherent structure so long as it remains a repair equilibrium: so long as the
constraint propagation reaching each node is sufficient to maintain that node
within the admissibility region of the tree's semantic manifold. When the
constraint propagation falls below a local threshold, the node crosses the
boundary of the admissibility region. Below the boundary, the system's own
dynamics do not restore it to admissibility; they carry it further out. This
is the irreversibility condition. And because semantic structures are
hierarchically organized, boundary failure at a peripheral node propagates
inward, undermining the constraint infrastructure for the nodes that depend
on the now-failed peripheral region.

The paper is organized as follows. Section~2 defines the formal apparatus:
semantic trees, admissibility weights, constraint propagation, and attenuation.
Section~3 establishes the repair-equilibrium interpretation of each node and
defines the key quantities $S(v)$, $\theta(v)$, and the repair ratio
$\rho(v) = S(v)/\theta(v)$. Section~4 proves the Semantic Stability Theorem
for individual nodes and extends it to branches via the branch coherence
condition $C(\mathcal{B})$. Section~5 analyzes boundary dynamics and proves
the propagation theorem: local boundary failure produces inward crack-like
failure fronts. Section~6 develops the RSVP field interpretation of the entire
framework. Section~7 discusses the linguistic and cognitive applications, with
particular attention to language death and institutional memory collapse as
cases of repair-threshold crossing. Section~8 concludes.

\section{Formal Apparatus}

\subsection{Semantic Trees}

\begin{definition}[Semantic tree]
A \textit{semantic tree} $\mathcal{T} = (V, E, \ell, \Phi_0)$ consists of
a finite directed tree $(V, E)$ with root $r \in V$, a labeling function
$\ell: V \to \Sigma$ assigning semantic content to each node from a semantic
domain $\Sigma$, and a root constraint $\Phi_0 \in \Phi$ specifying the
initial semantic density at the root.
\end{definition}

The edges in $E$ are directed away from the root: each edge $(u, v) \in E$
denotes that $u$ is the parent of $v$. The tree encodes how a root concept
or claim ramifies into a hierarchy of derived concepts, sub-claims, dependent
interpretations, or derived meanings. We write $\mathrm{Anc}(v)$ for the set
of ancestors of $v$, including $v$ itself, and $\mathrm{Desc}(v)$ for the
set of descendants, including $v$.

\begin{definition}[Admissibility weight]
An \textit{admissibility weight} $w: V \to [0, 1]$ assigns to each node $v$
a value $w(v)$ representing that node's capacity to transmit semantic
constraint to its descendants. A node with $w(v) = 1$ transmits constraint
without loss; a node with $w(v) = 0$ is a complete semantic barrier.
\end{definition}

Admissibility weights are distinct from the semantic content $\ell(v)$. A
node may have rich semantic content ($\ell(v)$ is complex) but low transmission
capacity ($w(v)$ is small), meaning it is itself coherent but fails to
constrain interpretation below it. This can occur when a concept is highly
ambiguous, when it connects poorly to the surrounding semantic network, or when
it has been separated from its constraining context.

\subsection{Alpha Attenuation}

\begin{definition}[Attenuation function]
An \textit{attenuation function} $\alpha: E \to (0, 1]$ assigns to each edge
$(u, v)$ a transmission coefficient $\alpha(u, v)$ representing the fraction
of semantic constraint that survives the traversal of that edge. Attenuation
is multiplicative along paths: for a path $r = v_0, v_1, \ldots, v_k = v$,
the path attenuation is
\[
\alpha(r \to v) = \prod_{i=0}^{k-1} \alpha(v_i, v_{i+1}).
\]
\end{definition}

Multiplicativity captures the essential feature of hierarchical semantic
degradation: constraint weakens with each additional level of derivation, and
the weakening is compounding. A chain of mostly-faithful transmissions
$\alpha = 0.9$ over ten steps delivers only $0.9^{10} \approx 0.35$ of the
original constraint at the tenth level --- just over a third of what the root
specified. Over twenty steps it delivers $0.9^{20} \approx 0.12$.

This is not an artifact of the formalism. It reflects the genuine difficulty
of maintaining interpretive fidelity across long chains of derivation, analogy,
translation, or institutional transmission.

\subsection{Constraint Propagation and Semantic Support}

\begin{definition}[Semantic support]
The \textit{semantic support} $S(v)$ at node $v$ is the total constraint
influence reaching $v$ from its ancestors, weighted by path attenuation and
ancestral admissibility:
\[
S(v) = \sum_{u \in \mathrm{Anc}(v)} \alpha(r \to u)\, w(u)\, K(u, v),
\]
where $K(u, v) \geq 0$ is a kernel measuring the direct relevance of ancestor
$u$'s content to descendant $v$'s admissibility, and $\alpha(r \to u)$ is
the path attenuation from root to $u$.
\end{definition}

The kernel $K(u, v)$ encodes the semantic architecture of the tree: how much
the meaning at $u$ constrains interpretation at $v$, independent of the
transmission efficiency. In the simplest case, $K(u, v) = 1$ for direct
ancestors and $0$ otherwise, reducing semantic support to pure path attenuation.
More realistic models allow $K$ to decay smoothly with semantic distance,
enabling off-path influence from conceptually close but structurally distant
ancestors.

In the uniform case with $K(u, v) = 1$ for all $u \in \mathrm{Anc}(v)$ and
constant $\alpha$ per edge, the semantic support becomes a sum of geometric
terms:
\[
S(v) = w(r) + \alpha\, w(v_1) + \alpha^2\, w(v_2) + \cdots + \alpha^{d(v)-1}\, w(v_{d(v)-1}),
\]
where $d(v)$ is the depth of $v$ and $v_0, v_1, \ldots, v_{d(v)-1}$ is the
path from root to $v$'s parent.

\section{Nodes as Repair Equilibria}

\subsection{The Admissibility Threshold}

\begin{definition}[Admissibility threshold]
Each node $v \in V$ has an \textit{admissibility threshold} $\theta(v) > 0$,
the minimum semantic support required for $v$ to maintain its semantic identity
--- to function as a coherent node that successfully constrains its own
descendants and participates in the tree's semantic manifold.
\end{definition}

The threshold $\theta(v)$ is not a property of $v$'s content alone but of
$v$'s role in the tree. A node that is a conceptual hub --- from which many
other nodes derive their interpretation --- requires higher semantic support
to maintain because the consequences of its failure are larger. A peripheral
leaf node requires lower support because its failure does not cascade.

\begin{definition}[Repair ratio]
The \textit{repair ratio} at node $v$ is
\[
\rho(v) = \frac{S(v)}{\theta(v)}.
\]
A node is \textit{semantically alive} if $\rho(v) \geq 1$ and
\textit{semantically failing} if $\rho(v) < 1$.
\end{definition}

This is the direct semantic analogue of the rate-matching condition
$r_R \geq r_D$ from the repair ontology. Semantic support $S(v)$ is the
repair pressure reaching the node. The threshold $\theta(v)$ is the minimum
repair rate required for the node to persist as a coherent element of the
semantic manifold. The node is a repair equilibrium exactly when $\rho(v) \geq 1$.

\begin{remark}
The repair-equilibrium interpretation resolves an apparent circularity in
semantic theory. Semantic content at $v$ depends on the content of $v$'s
ancestors, but what counts as valid ancestral content depends on the tree's
overall coherence. The repair-equilibrium framework dissolves this circle:
the dependence is not circular but dynamic. Each node's semantic identity
is continuously re-established by the repair pressure reaching it from above.
The identity is not a fixed property but an ongoing achievement.
\end{remark}

\subsection{The Admissibility Manifold for Semantic Trees}

In the language of the repair ontology, the admissibility manifold for a
semantic tree is the set of all weight-threshold configurations
$(w, \theta, \alpha)$ under which every node is semantically alive:
\[
\mathcal{A}(\mathcal{T}) = \{(w, \theta, \alpha) : \rho(v) \geq 1 \text{ for all } v \in V\}.
\]
The boundary $\partial\mathcal{A}(\mathcal{T})$ is the set of configurations
at which at least one node is exactly at its admissibility threshold. These
are the critical configurations: any increase in attenuation, any decrease in
admissibility weight, or any increase in a node's threshold will push the
system outside $\mathcal{A}$.

\section{The Semantic Stability Theorem}

\subsection{Node-Level Stability}

\begin{theorem}[Node stability]
A node $v$ remains semantically alive under perturbation $\delta\alpha$ of
the edge attenuation if and only if
\[
S(v) - \theta(v) > \left|\frac{\partial S(v)}{\partial \alpha}\right| \cdot |\delta\alpha|.
\]
The left side is the repair surplus at $v$; the right side is the rate at which
perturbation depletes that surplus. Nodes with large surplus tolerate large
attenuation perturbations; nodes near the threshold $\rho(v) \approx 1$
tolerate almost none.
\end{theorem}

\begin{proof}
Under perturbation $\delta\alpha$, the semantic support changes to
$S(v) + (\partial S(v)/\partial\alpha)\,\delta\alpha$. The node remains
semantically alive when $S(v) + (\partial S(v)/\partial\alpha)\,\delta\alpha \geq \theta(v)$,
i.e., when $S(v) - \theta(v) \geq -(\partial S(v)/\partial\alpha)\,\delta\alpha$.
Taking absolute values and worst-case sign gives the stated condition.
\end{proof}

\subsection{Branch Coherence}

Individual node stability is necessary but not sufficient for tree-level
semantic persistence. A branch may contain nodes with individually positive
repair ratios while the weakest node in the branch provides insufficient
support to nodes that depend on it. We need a branch-level condition.

\begin{definition}[Branch stability index and resilience]
For a connected subtree $\mathcal{B} \subseteq V$, define two quantities.
The \textit{branch stability index} is the minimum repair ratio across all
nodes in the branch:
\[
\rho_{\min}(\mathcal{B}) = \min_{v \in \mathcal{B}} \rho(v) = \min_{v \in \mathcal{B}} \frac{S(v)}{\theta(v)}.
\]
The \textit{branch resilience} is the mean repair ratio:
\[
C(\mathcal{B}) = \frac{1}{|\mathcal{B}|} \sum_{v \in \mathcal{B}} \rho(v).
\]
These are distinct quantities with distinct roles. The stability index
determines survival; the resilience measures capacity to absorb further
perturbation.
\end{definition}

\begin{remark}
The distinction matters. Consider a branch with repair ratios $(2.0, 2.0, 0.1)$. Its
resilience is $C(\mathcal{B}) = 1.37$, suggesting apparent health, but its
stability index is $\rho_{\min} = 0.1$, meaning one node is already dead
and the branch is not a repair equilibrium. Averages mask the most dangerous
structure in any hierarchical system: the weakest link. Since semantic trees
transmit constraint hierarchically, a failed intermediate node severs support
to all its descendants regardless of the health of sibling nodes.
\end{remark}

\begin{theorem}[Semantic Stability Theorem]
A connected semantic subtree $\mathcal{B}$ is a repair equilibrium if and
only if $\rho_{\min}(\mathcal{B}) \geq 1$, i.e., every node in the branch
satisfies $\rho(v) \geq 1$. The branch resilience $C(\mathcal{B})$ measures
the surplus available to absorb future attenuation increases: a branch with
$C(\mathcal{B}) \gg 1$ can tolerate large perturbations before any node
approaches its admissibility threshold.
\end{theorem}

\begin{proof}
($\Rightarrow$) If $\mathcal{B}$ is a repair equilibrium then by definition
$\rho(v) \geq 1$ for all $v \in \mathcal{B}$, so
$\rho_{\min}(\mathcal{B}) = \min_v \rho(v) \geq 1$.

($\Leftarrow$) If $\rho_{\min}(\mathcal{B}) \geq 1$ then $\rho(v) \geq 1$
for all $v \in \mathcal{B}$ by definition of the minimum, so every node is
semantically alive and $\mathcal{B}$ is a repair equilibrium.

($\neg$: instability) If $\rho_{\min}(\mathcal{B}) < 1$, there exists
$v^* \in \mathcal{B}$ with $S(v^*) < \theta(v^*)$. Since
$\sum_{v} S(v) < \sum_{v} \theta(v)$ whenever any term satisfies
$S(v) < \theta(v)$ (by direct arithmetic; no continuity argument is needed),
the total repair pressure is insufficient. Furthermore, $v^*$'s failure
reduces its effective admissibility weight, decreasing $S(u)$ for all
$u \in \mathrm{Desc}(v^*)$. This may push additional nodes below threshold.
The branch cannot self-repair without external support.
\end{proof}

\begin{remark}
The Semantic Stability Theorem is the exact discrete analogue of the
rate-matching condition $r_R \geq r_D$ from the repair ontology. The stability
index $\rho_{\min}$ plays the role of the ratio $r_R/r_D$ evaluated at the
worst-performing site. The resilience $C(\mathcal{B})$ plays the role of
mean repair surplus, which governs how much additional degradation the system
can absorb before a node crosses the boundary.
\end{remark}

\subsection{The Critical Attenuation Level}

\begin{definition}[Critical attenuation]
For a semantic tree $\mathcal{T}$, the \textit{critical attenuation level}
$\alpha_c$ is the minimum uniform edge attenuation coefficient for which
every node remains semantically alive, i.e., for which
$\rho_{\min}(V) \geq 1$:
\[
\alpha_c = \inf\{\alpha^* : \rho_{\min}(V) \geq 1 \text{ under uniform } \alpha = \alpha^*\}.
\]
\end{definition}

\begin{proposition}[Attenuation threshold]
For a tree of depth $D$ with uniform admissibility weights $w$ and uniform
thresholds $\theta$, the critical attenuation level satisfies
\[
\alpha_c = \left(\frac{\theta}{w \cdot H_D}\right)^{1/D},
\]
where $H_D = \sum_{k=0}^{D-1} \alpha^k$ is the geometric sum evaluated at
the boundary. Below $\alpha_c$, semantic support at the deepest nodes falls
below threshold and the tree begins collapsing from the leaves inward.
\end{proposition}

The key feature of this threshold is its non-linearity in depth. For deep
trees, small reductions in $\alpha$ below $\alpha_c$ produce large reductions
in $S(v)$ at the leaves, because attenuation compounds multiplicatively. This
explains why deep semantic structures --- long chains of derivation, complex
institutional memories, multi-generational linguistic traditions --- are
disproportionately vulnerable to small increases in transmission loss.

\section{Boundary Dynamics and Failure Propagation}

\subsection{Repair Intensity at Semantic Boundaries}

From the repair ontology, repair intensity is concentrated near the boundary
$\partial\mathcal{A}$:
\[
\iota(v) \propto \frac{1}{d(v,\, \partial\mathcal{A}(\mathcal{T}))} = \frac{1}{\rho(v) - 1}
\quad \text{for } \rho(v) > 1.
\]
Nodes with $\rho(v)$ close to $1$ are near the semantic boundary. They require
continuous and intense repair pressure to remain semantically alive. Nodes with
$\rho(v) \gg 1$ are deep inside the admissibility region and require relatively
little active maintenance.

\begin{theorem}[Critical Repair Divergence]
As a node approaches the admissibility boundary from above,
\[
\lim_{\rho(v) \to 1^+} \iota(v) = +\infty.
\]
That is, the maintenance cost required to keep a semantically marginal node alive
diverges as its repair ratio approaches the critical threshold. No finite
external intervention can sustain a node at exactly $\rho(v) = 1$ against any
additional perturbation.
\end{theorem}

\begin{proof}
From the definition $\iota(v) \propto 1/(\rho(v) - 1)$ for $\rho(v) > 1$,
the limit $\rho(v) \to 1^+$ sends $\rho(v) - 1 \to 0^+$, and therefore
$\iota(v) \to +\infty$. For the second claim: maintaining $\rho(v) = 1$
against a perturbation $\delta > 0$ to the degradation rate requires increasing
$S(v)$ by $\delta \cdot \theta(v)$, which demands a repair flow increase
proportional to $1/(\rho(v)-1) \cdot \delta$. As $\rho(v) \to 1^+$ this
required flow diverges for any fixed $\delta > 0$.
\end{proof}

The Critical Repair Divergence theorem gives the formalism its most directly
observable consequence. It explains why late-stage interventions in failing
systems are disproportionately expensive: an endangered language requiring its
last hundred speakers to sustain it demands more repair work per node than
a thriving language with millions; a dying institution requires increasingly
heroic maintenance efforts from a shrinking core; a marginal scientific paradigm
requires more and more defensive effort to maintain coherence against accumulating
anomalies. The cost does not increase linearly as the boundary approaches. It
diverges. This is the formal mechanism behind the empirical observation that
preventive maintenance is always cheaper than corrective maintenance: the
corrective regime operates near the boundary, where $\iota$ is large; the
preventive regime operates in the interior, where $\iota$ is small.

The divergence also implies a critical choice structure near the boundary. A
system cannot remain indefinitely near $\rho(v) = 1$; the maintenance cost
required to sustain it there is unbounded. Any real system with finite repair
capacity must either recover to $\rho(v) \gg 1$ (by external input of support)
or cross the boundary to $\rho(v) < 1$ (entering irreversible attenuation).
The boundary is not a stable equilibrium; it is a saddle point between recovery
and collapse.

In practice, the boundary-concentrated nodes are the conceptually peripheral
nodes of the tree: those at high depth, those with many competing interpretive
pressures, those that depend on long chains of derivation from the root, and
those whose admissibility weights are locally low due to contextual
disconnection. These are precisely the nodes most vulnerable to attenuation
increase and most likely to fail first.

\subsection{Failure Front Propagation}

The following theorem establishes that boundary failure does not remain
localized. It propagates inward.

\begin{theorem}[Semantic crack propagation]
Let $v^* \in \partial\mathcal{A}(\mathcal{T})$ be a boundary node with
$\rho(v^*) = 1$ that undergoes a perturbation pushing it below threshold
($\rho(v^*) < 1$). Then for every descendant $u \in \mathrm{Desc}(v^*)$,
the semantic support satisfies
\[
S(u) \leq S_{\mathrm{pre}}(u) - \Delta(v^*, u),
\]
where $\Delta(v^*, u) > 0$ is the loss of ancestral support due to
$v^*$'s failure. In particular, there exists a critical depth
$d_c$ below $v^*$ such that all nodes at depth $\geq d_c$ below $v^*$
will fall below their thresholds following $v^*$'s failure, regardless of
their pre-failure repair ratios.
\end{theorem}

\begin{proof}
When $v^*$ crosses below threshold, its admissibility weight $w(v^*)$
effectively drops, reducing the semantic support it contributes to all
descendants. For $u \in \mathrm{Desc}(v^*)$, the contribution of $v^*$
to $S(u)$ is $\alpha(r \to v^*)\, w(v^*)\, K(v^*, u)$. The failure of $v^*$
reduces $w(v^*)$ by some $\delta w > 0$, reducing $S(u)$ by
$\Delta(v^*, u) = \alpha(r \to v^*)\, \delta w\, K(v^*, u) > 0$.
Nodes sufficiently deep below $v^*$ have their repair surplus
$S(u) - \theta(u)$ already depleted by cumulative attenuation. The reduction
$\Delta(v^*, u)$ is sufficient to push these nodes below threshold. The
critical depth $d_c$ is where $\Delta(v^*, u) \geq S_{\mathrm{pre}}(u) - \theta(u)$
for all nodes at that depth.
\end{proof}

The analogy to materials failure is precise. In fracture mechanics, a surface
crack propagates inward when the stress intensity factor at the crack tip
exceeds the material's fracture toughness. Here, semantic failure propagates
inward from the boundary when the reduction in support at a failed node exceeds
the repair surplus of the nodes it was supporting. The structure appears stable
until the boundary defect exceeds a critical magnitude; then the failure front
advances into the interior.

\begin{corollary}[Extinction as repair failure]
The total semantic collapse of a tree $\mathcal{T}$ --- the state in which
no node remains semantically alive --- is the endpoint of failure front
propagation from all boundary nodes simultaneously falling below threshold.
This is not information loss: the semantic content $\ell(v)$ may remain
stored somewhere. It is repair failure: the constraint propagation network
that maintained the admissibility of each node has disintegrated faster than
any sub-network could compensate.
\end{corollary}

This corollary unifies several apparently distinct phenomena under a single
formal description. Language death is total semantic collapse of a linguistic
tree. Species extinction is total semantic collapse of an ecological constraint
network. Institutional collapse is total semantic collapse of a normative and
procedural tree. Memory loss is total semantic collapse of a reconstructive
constraint network. In each case the mechanism is the same: boundary failure
propagating inward faster than internal repair can compensate.

\section{The RSVP Field Interpretation}

The RSVP field triple $(\Phi, v, S)$ provides a field-theoretic setting for
the semantic stability framework. The correspondence is not metaphorical. We
establish explicit identifications that make the semantic tree formalism a
special case of RSVP field dynamics on a discrete graph substrate.

\subsection{Explicit Field Identifications}

\begin{definition}[RSVP-semantic correspondence]
Given a semantic tree $\mathcal{T} = (V, E, \ell, \Phi_0)$ with admissibility
weights $w$, attenuation $\alpha$, and repair ratios $\rho$, the RSVP fields
are identified as follows:
\begin{align}
\Phi(v) &= S(v) \label{eq:phi} \\
|v_{ij}| &= \alpha(i,j) \quad \text{for each edge } (i,j) \in E \label{eq:vecfield} \\
S_{\mathrm{RSVP}}(v) &= -\log \rho(v) \label{eq:entropy}
\end{align}
where $\Phi(v)$ is the scalar field at node $v$, $|v_{ij}|$ is the magnitude
of the vector field on edge $(i,j)$, and $S_{\mathrm{RSVP}}(v)$ is the
entropy field at node $v$.
\end{definition}

These identifications are not free choices. Each is forced by the requirement
that the RSVP stability condition coincide with the Semantic Stability Theorem.

Equation~(\ref{eq:phi}): The scalar field $\Phi$ encodes semantic density in
the RSVP framework. Identifying $\Phi(v) = S(v)$ means that semantic support
--- the accumulated constraint propagation reaching a node --- is exactly the
local density of the scalar field. High $\Phi$ at a node means high constraint
support; the root, which has maximal constraint specification, sits at the
global maximum of $\Phi$. Peripheral nodes at high depth, where attenuation
has compounded, have low $\Phi$, consistent with their low semantic support.

Equation~(\ref{eq:vecfield}): The vector field $v$ encodes constraint
propagation velocity in the RSVP framework. Identifying its magnitude with
the edge attenuation coefficient makes the vector field a direct measure of
transmission fidelity. An edge with $\alpha(i,j) = 1$ is a perfect constraint
channel; $|v_{ij}| = 1$ means full propagation speed with no attenuation loss.
An edge with $\alpha(i,j) \ll 1$ is a near-barrier; $|v_{ij}| \ll 1$ means
the constraint flow is nearly stopped at that edge. Failure of the vector
field at an edge is directly attenuation failure.

Equation~(\ref{eq:entropy}): The identification $S_{\mathrm{RSVP}}(v) = -\log\rho(v)$
is the most consequential. It has several immediately verifiable properties.
When $\rho(v) = 1$, the node is exactly at the admissibility boundary and
$S_{\mathrm{RSVP}}(v) = 0$: zero entropy means the interpretation volume is
exactly at the threshold of admissibility. When $\rho(v) > 1$, the node is
inside the admissibility manifold and $S_{\mathrm{RSVP}}(v) < 0$: negative
entropy corresponds to the constraint surplus, the degree to which the node's
interpretation is over-determined relative to its threshold. When $\rho(v) < 1$,
the node is outside the admissibility manifold and $S_{\mathrm{RSVP}}(v) > 0$:
positive entropy means the accessible interpretation volume has expanded beyond
what the constraint infrastructure can bound. Failure corresponds exactly to
positive entropy.

\begin{proposition}[RSVP stability condition]
Under the identifications above, the RSVP stability condition
\[
S_{\mathrm{RSVP}}(v) \leq 0 \quad \text{for all } v \in V
\]
is equivalent to the Semantic Stability Theorem's condition
$\rho_{\min}(V) \geq 1$.
\end{proposition}

\begin{proof}
$S_{\mathrm{RSVP}}(v) = -\log\rho(v) \leq 0$ iff $-\log\rho(v) \leq 0$ iff
$\log\rho(v) \geq 0$ iff $\rho(v) \geq 1$.
Taking the condition over all $v \in V$ gives $\rho_{\min}(V) \geq 1$.
\end{proof}

\begin{corollary}[Critical Repair Divergence in RSVP]
The Critical Repair Divergence theorem takes the following form in the RSVP
field language. As $S_{\mathrm{RSVP}}(v) \to 0^-$ (the node approaches the
admissibility boundary from inside),
\[
\iota(v) = \frac{1}{e^{-S_{\mathrm{RSVP}}(v)} - 1} \to +\infty.
\]
The repair intensity diverges as the entropy field approaches zero from below.
\end{corollary}

\subsection{Alpha Attenuation as Boundary Crossing}

The identification $|v_{ij}| = \alpha(i,j)$ makes the mechanism of boundary
crossing precise. Increasing attenuation on edge $(i,j)$ decreases $|v_{ij}|$,
reducing the constraint flow to $j$ and all descendants of $j$. By
equation~(\ref{eq:phi}), this decreases $\Phi$ at those nodes. By
equation~(\ref{eq:entropy}), this increases $S_{\mathrm{RSVP}}$. When
$S_{\mathrm{RSVP}}(v)$ crosses zero, the node exits the admissibility manifold:
\[
\alpha(i,j) \text{ falls below critical}
\;\Rightarrow\; S_{\mathrm{RSVP}}(j) > 0
\;\Rightarrow\; \rho(j) < 1
\;\Rightarrow\; j \notin \mathcal{A}(\mathcal{T}).
\]

Alpha attenuation is therefore not information loss in any Shannon sense.
Shannon entropy is a property of source distributions; it is invariant to
what the signal means. What attenuation destroys here is $S_{\mathrm{RSVP}}$
--- the constraint-entropy field that measures how tightly the admissibility
manifold bounds interpretation. The semantic content $\ell(v)$ may survive
perfectly while $S_{\mathrm{RSVP}}(v)$ crosses zero: the node is readable
but no longer repaired. This is the formal version of the stored-structure
versus living-equilibrium distinction. A preserved node has $\ell(v)$ intact
but $S_{\mathrm{RSVP}}(v) > 0$. A living node has $S_{\mathrm{RSVP}}(v) \leq 0$.

\subsection{Lamphrodyne Relaxation as Semantic Repair}

The lamphrodyne relaxation dynamics of RSVP smooth extreme local gradients in
$\Phi$ through the coupled dynamics of $v$ and $S_{\mathrm{RSVP}}$. Under the
semantic identification, lamphrodyne relaxation corresponds to the repair of
boundary nodes: when a peripheral node approaches $S_{\mathrm{RSVP}}(v) \to 0^-$,
the relaxation mobilizes constraint flow (increases $|v_{ij}|$ on incoming
edges), which increases $\Phi(v) = S(v)$ and decreases $S_{\mathrm{RSVP}}(v)$,
pushing the node back toward the interior.

When the lamphrodyne relaxation rate is insufficient to prevent $S_{\mathrm{RSVP}}$
from crossing zero, the node undergoes irreversible attenuation. The Spherepop
COLLAPSE outcome formalizes this at the computational level: once a node has
committed to a reduced semantic state, its bubble closes and $S_{\mathrm{RSVP}}$
for the collapsed region is reset to the new, lower-constraint manifold. The
repaired tree is not the original tree; it is a new repair equilibrium with
a different admissibility manifold, excluding the collapsed region.

\section{Linguistic and Cognitive Applications}

\subsection{Language Death as Threshold Crossing}

The repair-equilibrium framework gives a formal account of language death that
distinguishes it clearly from language change.

Language change occurs when a semantic tree undergoes repair that reaches the
admissibility manifold at a different point: the repaired language satisfies
the same constraints as the pre-change language but with different realizations.
Grammar changes, sound changes, and lexical innovations are all instances of
repair reaching new points on the manifold while maintaining the tree's overall
coherence. The branch coherence $C(V) \geq 1$ is preserved.

Language death occurs when $C(V)$ falls below $1$. This happens when the
transmission infrastructure --- the intergenerational and community-level
repair processes that sustain the language's semantic manifold --- is disrupted
faster than internal repair can compensate. The trigger is typically the
reduction of the speech community below a critical size, the interruption of
intergenerational transmission, or the rapid replacement of the language's
contextual infrastructure by a dominant competing language.

In the boundary-dynamics framework, language death proceeds by failure front
propagation. The peripheral nodes of the semantic tree --- rare vocabulary,
specialized registers, culturally specific pragmatic conventions, low-frequency
syntactic constructions --- are the first to lose their repair support. As the
speech community contracts or shifts, these nodes cross the admissibility
boundary first. Their failure reduces the support available to the intermediate
nodes that depended on them for contrast and contextual specificity. The
failure front advances inward toward the core vocabulary and basic grammar.
At the point of total collapse, even the core nodes --- the high-frequency,
low-attenuation, maximally supported nodes --- fall below threshold because
the entire constraint propagation network has disintegrated.

A critical implication: language death is not reversible by reintroducing the
semantic content. A dictionary of a dead language contains the nodes $\ell(v)$
but not the constraint infrastructure $S(v) \geq \theta(v)$ that made each
node semantically alive. Revitalization requires reconstructing the repair
infrastructure --- the speech community, the use contexts, the intergenerational
transmission mechanisms --- not merely the content.

\subsection{Institutional Memory Collapse}

Institutional memory fails through the same mechanism, formalized identically.

An institution's normative and procedural knowledge is a semantic tree: the
foundational principles of the institution are the root; the specific procedures,
exceptions, informal norms, and contextual knowledge are the peripheral nodes.
The admissibility weight of each node is the degree to which that piece of
knowledge is actively practiced, transmitted to new members, and reinforced by
the institution's ongoing activity. The threshold $\theta(v)$ of each node
is the minimum active-practice rate for that knowledge to remain functionally
accessible.

Personnel attrition is attenuation: as individuals who possess peripheral
procedural knowledge leave, the transmission coefficient $\alpha$ for that
knowledge falls. Initially this affects only the peripheral nodes. But if
the attrition rate exceeds the institution's onboarding and knowledge-transfer
capacity, the repair rate falls below the degradation rate and the branch
coherence of the institutional knowledge tree begins to fall. The failure front
propagates inward from specialized tacit knowledge toward the core operating
procedures.

The formal threshold is the same: $C(\mathcal{B}) \geq 1$ for knowledge
survival, $C(\mathcal{B}) < 1$ for irreversible loss. An institution that
re-hires former personnel with the intention of recovering lost knowledge
discovers that the nodes $\ell(v)$ may be recoverable --- the individual
remembers --- but the constraint infrastructure that made the knowledge
institutionally alive has collapsed. The individual's knowledge is no longer
repaired by the institution's ongoing practice; it is an isolated node
disconnected from the manifold.

\subsection{Cognitive Depth and the Gesture-Before-Symbol Relationship}

The gesture-before-symbol perspective developed in prior work holds that meaning
is prior to its linguistic encoding. In the repair-equilibrium framework, this
becomes a claim about the tree structure: gesture and embodied context are the
root nodes of the semantic tree, and linguistic symbols are intermediate and
peripheral nodes that depend on the root's constraint propagation for their
admissibility.

This ordering has a precise consequence: the peripheral linguistic nodes can
fail without causing root-level semantic failure, but root-level failure
causes total collapse. A community that loses a specific word (peripheral
node) retains the meaning (root constraint) and can repair the loss through
alternative expression. A community that loses the embodied, contextual,
gestural grounding of a meaning domain loses the constraint infrastructure for
the entire branch of symbols that depended on it. The symbols may survive in
dictionaries and grammars while being semantically dead in the repair-equilibrium
sense: no longer repaired by the constraint propagation from their gestural and
contextual roots.

This also explains the particular vulnerability of abstract terminology. Abstract
terms are high-depth nodes in the semantic tree: they depend on long chains of
derivation from more concrete roots, and each link in the chain contributes
multiplicative attenuation. The critical attenuation level $\alpha_c$ for a
deep abstract term is much higher than for a concrete high-frequency term,
meaning that abstract terminology requires substantially more active maintenance
to remain semantically alive. This is why technical vocabulary deteriorates
rapidly when a community stops actively practicing the domain it describes.

\section{Conclusion}

The Semantic Stability Theorem establishes that a semantic tree is a repair
equilibrium in precisely the formal sense developed in the companion paper on
repair ontology. Each node is semantically alive when its repair ratio
$\rho(v) = S(v)/\theta(v) \geq 1$. The branch stability index
$\rho_{\min}(\mathcal{B})$ is the correct survival condition; the resilience
$C(\mathcal{B})$ measures the system's remaining tolerance for further
degradation. The Critical Repair Divergence theorem establishes that maintenance
cost diverges as any node approaches the admissibility boundary, explaining the
empirical universality of the observation that late-stage interventions are
disproportionately expensive. And the explicit RSVP identifications --- $\Phi = S$,
$|v_{ij}| = \alpha(i,j)$, $S_{\mathrm{RSVP}} = -\log\rho$ --- ground the
tree formalism in the field dynamics of the RSVP framework, making the
semantic stability condition the exact discrete analogue of the RSVP
field-stability condition $S_{\mathrm{RSVP}} \leq 0$.

The deepest consequence of the framework is a formal distinction between two
modes of persistence that ordinary language conflates.

\begin{definition}[Stored structure versus living equilibrium]
A \textit{stored structure} is a configuration in which the semantic content
$\ell(v)$ is preserved but $S_{\mathrm{RSVP}}(v) > 0$ for some or all nodes:
the constraint infrastructure has collapsed, and the content exists without
the repair dynamics that once kept it semantically alive.
A \textit{living equilibrium} is a configuration in which $\rho_{\min}(V) \geq 1$:
every node is actively maintained within the admissibility manifold by
ongoing constraint propagation.
\end{definition}

These are not different degrees of the same thing. They are categorically
different ontological states. A dead language dictionary is a stored structure:
the nodes $\ell(v)$ are intact, but $S_{\mathrm{RSVP}}(v) > 0$ everywhere ---
no repair dynamics sustain the admissibility of any entry. An abandoned
institution's policy manual is a stored structure. A fossil is a stored
structure. An archived codebase whose community of practice has dissolved is
a stored structure. A trained language model whose training culture no longer
exists is a stored structure: the interpolated topology of the residue, learned
from an immense graveyard of formerly living equilibria, performs a kind of
reconstruction over preserved content rather than maintaining live constraints.

In each case the content survives the death of the repair equilibrium.
Reconstruction from stored structure is possible but is not the same operation
as maintaining a living equilibrium. Revitalization of a dead language requires
reconstructing the repair infrastructure --- the speech community, the use
contexts, the intergenerational transmission --- not merely the content.
Reconstruction from the content alone produces an artifact, not a repair
equilibrium: a learned performance of a dead system, not the system itself.

Semantic stability is not a property of content. It is a property of
maintenance. What exists semantically are living repair equilibria. Everything
else is stored structure on a trajectory toward dissolution --- the accumulation
of successful repair while the repair lasted, and the residue of failed repair
thereafter.

\bigskip

\noindent\rule{0.4\textwidth}{0.4pt}

\vspace{0.5em}
\noindent{\small Flyxion / Independent Researcher. 2026.}

\end{document}

